% ============================================================
% ChatGIT: Granularity-Adaptive Conversational Repository Intelligence
% Springer LNCS
% Authors: Anish Gupta, Prakhar Sethi, Ritwik Bhattacharya,
%          Dr. Sonia Khetarpaul
% Institution: SNU — School of Engineering and Applied Science
% ============================================================

\documentclass[runningheads]{llncs}

\usepackage[T1]{fontenc}
\usepackage[utf8]{inputenc}
\usepackage{graphicx}
\usepackage{booktabs}
\usepackage{amsmath}
\usepackage{amssymb}
\usepackage{hyperref}
\usepackage{multirow}
\usepackage{array}
\usepackage{xspace}
\usepackage[inline]{enumitem}
\usepackage{microtype}
\usepackage[table]{xcolor}

\usepackage{tikz}
\usetikzlibrary{arrows.meta, positioning, backgrounds, calc, fit}

% ── Colours ───────────────────────────────────────────────
\definecolor{cBlue}  {HTML}{3B82F6}
\definecolor{cIndigo}{HTML}{6366F1}
\definecolor{cPurple}{HTML}{A855F7}
\definecolor{cGreen} {HTML}{16A34A}
\definecolor{cDark}  {HTML}{1E293B}

% ── Macros ────────────────────────────────────────────────
\newcommand{\chatgit}{\textsc{ChatGIT}\xspace}

% ══════════════════════════════════════════════════════════
\begin{document}

\title{ChatGIT: Granularity-Adaptive Conversational Repository
Intelligence with Multi-Turn Session Memory, Graph-Structural
Ranking, and Bidirectional Call-Context Augmentation}

\titlerunning{ChatGIT: Conversational Repository Intelligence}

\author{Anish Gupta\inst{1} \and
        Prakhar Sethi\inst{1} \and
        Ritwik Bhattacharya\inst{1} \and
        Sonia Khetarpaul\inst{1}}

\authorrunning{Gupta, Sethi, Bhattacharya, Khetarpaul}

\institute{
  Computer Science and Engineering,
  School of Engineering and Applied Science,
  Shiv Nadar University (SNU), India\\
  \email{\{ag801, ps385, rb213\}@snu.edu.in,\;
         sonia.khetarpaul@snu.edu.in}
}

\maketitle

% ══════════════════════════════════════════════════════════
\begin{abstract}
Navigating a large, unfamiliar software repository is a
time-consuming task that keyword search and IDE browsing handle
poorly. We present \chatgit, a Retrieval-Augmented Generation
(RAG) system that enables multi-turn conversational
question-answering over GitHub repositories, surfacing relevant
code with precise file and line citations. Beginning from a
keyword-search baseline (P@5\,=\,0.48, F1\,=\,0.60), we design
a pipeline that integrates AST-aware chunking, graph-structural
retrieval signals (PageRank, HITS hub/authority scoring, and
git-volatility weighting), intent-driven granularity adaptation,
multi-turn session memory with coreference resolution, and
bidirectional call-context augmentation. On a 200-query in-house
evaluation suite the system achieves \textbf{F1\,=\,0.88,
P@5\,=\,0.78, Recall@10\,=\,0.86}. On an 18-conversation
multi-turn evaluation set it reaches \textbf{MRR\,=\,0.389}
($3.7\times$ BM25, $p{=}0.013$) with \textbf{zero cross-turn
redundancy} versus 12.2\% for a vanilla RAG baseline.

\keywords{Conversational Code QA \and RAG \and Repository
Intelligence \and Session Memory \and Call Graph \and
Intent Classification \and PageRank \and Software Engineering}
\end{abstract}

% ══════════════════════════════════════════════════════════
\section{Introduction}
\label{sec:intro}

Modern software projects routinely span hundreds of files, several
programming languages, nested module hierarchies, and years of
commit history. When a developer joins a project mid-sprint, or
needs to audit an unfamiliar library, the standard tools ---
\texttt{grep}, IDE symbol search, or GitHub's own code search ---
require the user to already know what they are looking for.
They return raw matches, not explanations; they have no memory of
what was asked moments ago; and they treat every file as equally
important regardless of architectural role or change frequency.

Large-language-model (LLM) based coding assistants partially
address the comprehension gap, but three fundamental problems
remain unresolved in current tools:

\textbf{Cross-turn redundancy.}  Without retrieval-level session
state, every follow-up question re-retrieves the same
highest-scoring entry-point functions, flooding the context
window with repetitive content and forcing the model to
re-explain what it just explained.

\textbf{Structurally-blind retrieval.}  Cosine similarity over
embeddings is content-aware but structurally unaware.  It cannot
distinguish a core data-model class --- imported by forty other
files and therefore architecturally authoritative --- from a
one-off utility used nowhere else.  It also cannot tell whether a
file has been stable for two years or rewritten twice last week,
even though stability is a meaningful signal for which definition
is trustworthy.

\textbf{Fixed retrieval granularity.}  ``Where is
\texttt{authenticate()} defined?'' requires a single tight
function chunk.  ``Explain the authentication module'' requires
several file-level summaries.  A fixed \texttt{top\_k} and
uniform chunk size serves neither query type well.

\chatgit addresses all three limitations through a pipeline of
complementary technical components.  The contributions of this
paper are:

\begin{enumerate}[label=(\arabic*),leftmargin=*]
  \item A full conversational code QA pipeline
        (\S\ref{sec:architecture}--\S\ref{sec:components})
        combining semantic AST chunking, graph-structural
        retrieval signals, session-aware multi-turn memory,
        intent-adaptive granularity control, and post-generation
        line-number citation injection.
  \item An empirical evaluation (\S\ref{sec:eval}--\S\ref{sec:results})
        demonstrating a $+47\%$ F1 improvement over a
        keyword-search baseline and zero cross-turn redundancy in
        multi-turn conversations.
  \item Open-source release of the system, evaluation scripts,
        and annotated query sets.
\end{enumerate}

% ══════════════════════════════════════════════════════════
\section{Background and Related Work}
\label{sec:related}

\subsection{Retrieval-Augmented Generation}

RAG~\cite{lewis2021rag} grounds LLM responses by retrieving
relevant passages from an external corpus before generation,
reducing hallucination and enabling knowledge-intensive tasks.
Standard RAG pipelines use a vector store indexed with dense
embeddings and retrieve by cosine similarity~\cite{guo2024lightrag}.
Recent surveys~\cite{petroni2024irrag,abootorabi2025multimodal}
identify multi-hop reasoning and long-document grounding as open
challenges; code repositories present both simultaneously: a
question about one function may require evidence from its callers,
its callees, and the module-level description.

\subsection{Code Representation and Retrieval}

BM25~\cite{robertson2009bm25} provides a strong lexical baseline
for code search and remains competitive on identifier-heavy
queries.  CodeBERT~\cite{feng2020codebert} and
GraphCodeBERT~\cite{guo2021graphcodebert} extend BERT
pre-training to program tokens and data-flow graphs respectively,
producing dense representations that capture semantic similarity
across variable renames and minor refactors.
RepoCoder~\cite{zhang2023repocoder} proposes iterative
retrieval-generation for code completion, re-issuing retrieval
queries after each generation step.
RepoAgent~\cite{liu2024repoagent} uses an LLM with call-graph
traversal to produce repository-level documentation.  None of
these systems address multi-turn conversational QA, session
coherence, or retrieval granularity adaptation.

\subsection{Graph-Based Ranking}

PageRank~\cite{brin1998pagerank} assigns importance scores to
nodes in a directed graph by propagating rank through edges;
it has been applied to software dependency graphs for bug
prediction and change-impact analysis.
HITS~\cite{kleinberg1999hits} decomposes node importance into
two complementary scores: \emph{hub} (a node that links to many
good authorities) and \emph{authority} (a node cited by many
good hubs).  Bavota et al.~\cite{bavota2015changeadvisor} apply
link-analysis techniques to software co-change graphs.
Graph Attention Networks~\cite{velickovic2018gat} provide a
learned alternative to hand-crafted graph centrality.  We combine
PageRank and HITS as lightweight, interpretable structural signals
within our retrieval rescoring stage.

\subsection{Multi-Turn and Conversational QA}

CoQA~\cite{reddy2019coqa} and QuAC~\cite{choi2018quac}
established benchmarks for multi-turn QA over flat text passages.
Neither addresses code repositories, call-graph structure, or
retrieval granularity.  Coreference resolution in conversational
QA --- determining what ``it'' or ``the function I mentioned
earlier'' refers to --- has been studied in open-domain dialogue
but not applied to code-specific retrieval pipelines.

\subsection{Commercial AI Coding Assistants}

GitHub Copilot Chat~\cite{githubcopilot},
Sourcegraph Cody~\cite{cody}, Cursor~\cite{cursor},
Aider~\cite{aider}, and Continue.dev~\cite{continuedev} all
offer conversational interfaces over code.  All maintain
conversation history within a session, and several perform
some form of codebase indexing (Cody uses Sourcegraph's code
graph; Cursor uses AST-based chunk embeddings).
What none of them expose is \emph{retrieval-level} session state
--- the ability to suppress chunks that were already retrieved
in prior turns, boost files that are topically active, and
expand ambiguous follow-up queries using the session's function
discussion history.  Intent-adaptive retrieval parameter tuning
(adjusting $k$, rerank depth, and chunk-type preference per
query intent) is likewise absent from all surveyed tools.

% ══════════════════════════════════════════════════════════
\section{System Architecture}
\label{sec:architecture}

\chatgit is organised into four pipeline stages
(Figure~\ref{fig:pipeline}).

\textbf{Stage A: Offline Indexing.}
Performed once per repository.  The repository is cloned, source
files are parsed with language-specific AST analysers, and the
resulting function and class nodes are converted into token-aware
chunks.  Chunks are embedded and stored in a ChromaDB vector
index.  Concurrently, a function call graph is built with
NetworkX, and git commit history is analysed to derive file
volatility scores.

\textbf{Stage B: Query Pre-Processing.}
For each user turn, the raw query is (i) classified into one of
four intent categories (LOCATE, EXPLAIN, SUMMARIZE, DEBUG),
which immediately reconfigures retrieval parameters; (ii) expanded
via coreference resolution using the session's history; and
(iii) annotated with a session summary that is later prepended
to the LLM prompt.

\textbf{Stage C: Retrieval and Rescoring.}
The expanded query is issued against the vector index.  Retrieved
candidates are rescored by a chain of structural signals:
hybrid PageRank + query-conditioned attention, HITS hub/authority
scores, git-volatility weights, and session-memory adjustments.
A cross-encoder reranker with a per-file diversity cap produces
the final ranked list, which is augmented with bidirectional
call-context neighbours.

\textbf{Stage D: Generation.}
The selected chunks, session summary, and caller/callee context
are assembled into an intent-budget prompt and passed to the LLM.
The generated response is post-processed to inject precise
\texttt{file:line} citations, and the session state is updated
for the next turn.

%% ── PIPELINE FIGURE ─────────────────────────────────────
\begin{figure}[t]
\centering
\begin{tikzpicture}[
  %% shared styles
  sbox/.style 2 args={
    draw=#1!55!black, fill=#1!8, rounded corners=4pt,
    line width=0.75pt, font=\footnotesize,
    text width=#2, align=left, inner sep=7pt,
  },
  arr/.style={-{Stealth[length=5pt,width=4pt]},
              line width=0.9pt, color=#1},
  darr/.style={-{Stealth[length=4pt,width=3pt]},
               line width=0.65pt, dashed, color=#1},
  lbl/.style={font=\scriptsize, text=#1!70!black},
]

%%── Row 1 ──────────────────────────────────────────────
\node[sbox={cBlue}{5.3cm}] (A) at (0,0) {%
  \textbf{\textcolor{cBlue!80!black}{A\enspace Offline Indexing}}\\[4pt]
  \textbullet~Clone repository (GitHub URL)\\
  \textbullet~AST parse: function/class nodes\\
  \textbullet~Token-aware chunking + module summaries\\
  \textbullet~BGE-small-en-v1.5 embeddings\\
  \textbullet~ChromaDB vector index\\
  \textbullet~Build call graph (NetworkX)\\
  \textbullet~Git volatility analysis\\[2pt]
  \textit{one-time per repository}%
};

\node[sbox={cIndigo}{5.3cm}, right=0.6cm of A] (B) {%
  \textbf{\textcolor{cIndigo!80!black}{B\enspace Query Pre-Processing}}\\[4pt]
  \textbullet~Receive user query\\
  \textbullet~Intent classification:\\
  \hspace*{8pt}LOCATE / EXPLAIN / SUMMARIZE / DEBUG\\
  \textbullet~Set $k$, rerank depth, context budget\\
  \textbullet~Coreference resolution (3 strategies)\\
  \textbullet~Prepend session history summary\\[2pt]
  \textit{online, per query}%
};

%%── Row 2 ──────────────────────────────────────────────
\node[sbox={cPurple}{5.3cm}, below=0.7cm of A] (C) {%
  \textbf{\textcolor{cPurple!80!black}{C\enspace Retrieval \& Rescoring}}\\[4pt]
  \textbullet~Cosine-similarity vector retrieval\\
  \textbullet~Hybrid PageRank + QC-attention score\\
  \textbullet~HITS hub/authority boost (by intent)\\
  \textbullet~Git-volatility weight\\
  \textbullet~Session-memory score adjustment\\
  \textbullet~Cross-encoder rerank (MiniLM-L-6)\\
  \textbullet~Bidirectional call-context augmentation%
};

\node[sbox={cGreen}{5.3cm}, right=0.6cm of C] (D) {%
  \textbf{\textcolor{cGreen!80!black}{D\enspace Generation}}\\[4pt]
  \textbullet~Assemble intent-budget prompt\\
  \textbullet~Session summary + caller/callee context\\
  \textbullet~Groq Llama-3.1-8b-instant inference\\
  \textbullet~Extract code identifiers from answer\\
  \textbullet~Inject file:line\textsubscript{start}--line\textsubscript{end} citations\\
  \textbullet~Update session memory state\\[2pt]
  \textit{online, per query}%
};

%%── Arrows ──────────────────────────────────────────────
%% User query enters B
\draw[arr=cIndigo]
  ([xshift=1.3cm]B.north) -- ++(0,0.5)
  node[above, lbl=cIndigo] {User Query};

%% A → C  (offline index feeds retrieval)
\draw[darr=cBlue] (A.south) -- (C.north)
  node[midway, right, lbl=cBlue] {index \&\\[-1pt]signals};

%% B → C  (query flows to retrieval — curved)
\draw[arr=cIndigo]
  (B.south) to[out=260, in=55] (C.north east);

%% C → D  (retrieval → generation)
\draw[arr=cPurple] (C.east) -- (D.west);

%% D → B  (session update, dashed back)
\draw[darr=cGreen]
  (D.north) -- ++(0,0.35) -| ([xshift=-0.6cm]B.north)
  node[near start, above, lbl=cGreen] {session update};

%% Answer out of D
\draw[arr=cGreen]
  (D.south) -- ++(0,-0.45)
  node[below, lbl=cGreen] {Answer + file:line citations};

\end{tikzpicture}
\caption{
  \chatgit pipeline. \textbf{Stage A} (blue) runs once when a
  repository is loaded. \textbf{Stages B--D} (indigo, purple,
  green) execute for every user query. Dashed arrows show offline
  signals feeding the live retrieval stage.
}
\label{fig:pipeline}
\end{figure}

% ══════════════════════════════════════════════════════════
\section{Technical Components}
\label{sec:components}

\subsection{Token-Aware AST Chunking and Module Summaries}

Language-specific AST parsers (Python \texttt{ast}; regex-based
extractors for JS/TS/Java/C++) identify function and class
boundaries.  Each boundary becomes the primary chunk unit,
respecting semantic coherence: a function is never split
mid-body.  Token budgets are enforced at MAX\,=\,512 tokens
with an overlap of 64 tokens and a maximum of 6 chunks per
top-level node.

In addition, a \emph{module-summary chunk} is synthesised for
each file by concatenating docstrings and exported symbol
signatures into a compact $\leq$200-token descriptor.  These
summaries are assigned a $1.50\times$ retrieval boost for
SUMMARIZE-intent queries, giving high-level ``explain the module''
questions a dedicated, correctly-scoped retrieval target.

This contrasts with sliding-window approaches (used by BM25
baselines and RepoCoder) that split functions at arbitrary token
boundaries, discarding structural context.

\subsection{Git-History Volatility Weighting}

Not all stable-looking code is equally trustworthy, and not all
actively-developed code is equally relevant.  We compute a
three-factor volatility score per file from up to 500 recent
commits:
\begin{equation}
\text{vol}(f) =
  0.5\,\text{freq}(f) + 0.3\,\text{recency}(f) + 0.2\,\text{diversity}(f)
\label{eq:vol}
\end{equation}
where \textit{freq} is the normalised commit count,
\textit{recency} is an exponentially-decayed days-since-last-change
score, and \textit{diversity} is the normalised count of distinct
committers.  A stable, foundational file receives a retrieval
weight of up to $1.40\times$ while a file in active flux is
weighted at $1.00\times$.

A co-change graph built from commit co-occurrence additionally
captures logical coupling between files that share no import
edges but are consistently modified together, surfacing hidden
dependencies invisible to the call graph.

\subsection{Hybrid PageRank and Query-Conditioned Graph Attention}

We combine two graph-structural signals:
\begin{equation}
h(v,q) = \alpha(q)\cdot\text{PR}(v)
        +(1{-}\alpha(q))\cdot\text{QCA}(v,q)
\label{eq:hybrid}
\end{equation}
$\text{PR}(v)$ is the standard PageRank score of node $v$ in the
file import graph.  $\text{QCA}(v,q)$ is a query-conditioned
neighbourhood attention score:
\[
\text{QCA}(v,q)
  = 0.6\,\text{sim}(e_v, e_q)
  + 0.4\,\textstyle\overline{\text{sim}}_{u\in\mathcal{N}(v)}(e_u, e_q)
\]
where $e_v$ is the node's embedding, $e_q$ is the query embedding,
and $\mathcal{N}(v)$ is the one-hop neighbourhood.
The mixture weight $\alpha(q)\in\{0.25, 0.50, 0.70\}$ adapts to
query specificity: narrow identifier queries weight local
embedding similarity ($\alpha$ low) while broad conceptual
questions weight global rank ($\alpha$ high).
The boost is applied multiplicatively:
$\text{score} \mathrel{\times}= 1 + 0.05\,h(v,q)$ when
$h(v,q) > 0.3$.

\subsection{HITS Hub/Authority Scoring}

HITS~\cite{kleinberg1999hits} separates node importance into
two complementary roles via mutual reinforcement:
\begin{equation}
  \text{hub}(v)
    = \textstyle\sum_{u:\,v\to u}\text{auth}(u),\qquad
  \text{auth}(v)
    = \textstyle\sum_{u:\,u\to v}\text{hub}(u)
\label{eq:hits}
\end{equation}
We run HITS on both the file import graph and the function call
graph.  \emph{Hub nodes} are orchestrators: entry-point scripts,
CLI commands, and app factories that import or delegate to many
downstream modules.  \emph{Authority nodes} are core utilities:
data models, base classes, and shared helpers that are imported
by many others.

We route these scores by intent.  For SUMMARIZE queries
(``explain the architecture''), hub scores are boosted because
the user needs entry-point, high-level context.  For LOCATE and
EXPLAIN queries (``where is this function?'', ``how does X
work?''), authority scores are boosted because the target is
typically a widely-used, well-defined unit.

\subsection{Multi-Turn Session-Aware Retrieval Memory}

A per-session \texttt{SessionRetrievalMemory} object maintains
state across turns and adjusts candidate scores before reranking:
\begin{equation}
s'(c) = s_0(c)\cdot r(c)\cdot z(c)
\label{eq:session}
\end{equation}
$r(c)$ is a redundancy penalty: 0.30 if the chunk was retrieved
this turn, 0.60 if one turn ago, 0.85 if two or more turns ago.
$z(c) = 1 + 0.30\,w_f$ is a zone-coherence bonus proportional
to the current weight of file $f$ in the session's active-file
set; file weights are initialised to 1.0 on first retrieval and
decay by $0.80^t$ per subsequent turn.  Functions explicitly
discussed in prior answers receive a $1.20\times$ boost when
they reappear as candidates.

Three coreference resolution strategies expand ambiguous queries
before retrieval:
\begin{enumerate}[label=(\roman*),leftmargin=*]
  \item \textbf{Temporal back-references}:
        ``the function I asked in the beginning'' is expanded
        with the verbatim first query text and the first
        discussed function name.
  \item \textbf{Repeat queries}:
        ``again?'' or ``more?'' are replaced with the previous
        query text.
  \item \textbf{Pronoun references}:
        short queries containing ``it'', ``this function'', etc.\
        are appended with the last discussed function name as a
        context hint.
\end{enumerate}
A rolling session summary (prior questions, recently discussed
functions, active files) is prepended to every LLM prompt,
keeping the model oriented across turns without re-transmitting
full conversation history.

\subsection{Intent-Driven Granularity-Adaptive Retrieval}

A keyword-pattern intent classifier assigns one of four intent
classes and simultaneously sets five retrieval parameters
(Table~\ref{tab:intent_config}).  The four intents reflect
qualitatively different information needs: LOCATE requires a
single precise chunk; EXPLAIN requires function-level context
with call neighbourhood; SUMMARIZE requires breadth over multiple
file summaries; DEBUG requires the largest $k$ to surface
call stacks and error paths.

\begin{table}[t]
\centering
\caption{Intent-specific retrieval configuration.}
\label{tab:intent_config}
\small
\begin{tabular}{lrrrrr}
\toprule
Intent & $k$ & Rerank $n$ & Max/file & Ctx budget & Chunk boost \\
\midrule
LOCATE    & 25 & 6  & 2 & 4\,K  & $1.40\times$ statement \\
EXPLAIN   & 20 & 8  & 3 & 6\,K  & $1.20\times$ function  \\
SUMMARIZE & 12 & 5  & 5 & 7\,K  & $1.50\times$ module    \\
DEBUG     & 30 & 10 & 4 & 6.5\,K& $1.10\times$ mixed     \\
\bottomrule
\end{tabular}
\end{table}

\subsection{Cross-Encoder Reranking with File-Diversity Capping}

After vector retrieval and graph rescoring, the top-$k$ candidates
are reranked by \texttt{ms-marco-MiniLM-L-6-v2}~\cite{reimers2019sentencebert},
a lightweight cross-encoder trained for passage relevance.
After reranking, a \emph{file-diversity cap} enforces the
\texttt{max\_per\_file} limit from Table~\ref{tab:intent_config},
preventing a single highly-similar file from consuming most
of the context budget when many of its chunks score similarly.

\subsection{Bidirectional Call-Context Neighbourhood Augmentation}

A retrieved function chunk without its calling context is
incomplete for explain and debug questions: understanding what
a function does requires knowing \emph{where it is called from}
and \emph{what it calls internally}.  For the top-5 seed chunks
with cosine similarity $>0.3$, we retrieve both direct callers
and direct callees from the call graph and augment the candidate
pool with a proportional boost:
\begin{equation}
\text{boost}(c_{\text{nb}})
  \mathrel{+}= 0.06 \times s_{\text{seed}}
\label{eq:callctx}
\end{equation}
Caller chunks are tagged \texttt{[CALLER]} and callee chunks are
tagged \texttt{[CALLEE]} in the prompt, giving the LLM explicit
structural provenance.  The component is active for EXPLAIN and
DEBUG intents only.  The proportional formulation in
Eq.~\eqref{eq:callctx} proved critical: an earlier fixed
$+0.05$ boost degraded MRR by $9.7\%$ by injecting
low-confidence neighbours at equal weight regardless of seed
quality.

\subsection{Post-Generation Line-Number Enhancement}

After generation, \texttt{ImprovedCodeSnippetExtractor}
post-processes the answer: code identifiers and function names
are extracted from the response text, located in the repository
index, and \texttt{file:line\_start--line\_end} annotations are
injected inline.  This step is deterministic and requires zero
additional LLM calls, providing navigable citations at no
added latency.

% ══════════════════════════════════════════════════════════
\section{Experimental Setup}
\label{sec:eval}

\subsection{Datasets}

\textbf{Setting A — in-house query suite (200 queries).}
We collected 200 natural-language queries manually authored by
the system developers, each mapped to one or more gold code
spans within five widely-used Python repositories: Flask,
Requests, FastAPI, Celery, and Click.  Queries cover
authentication flows, database operations, initialisation
sequences, configuration parsing, CLI routing, and error
handling.

\textbf{Setting B — multi-turn evaluation set (18 conversations).}
We curated 18 multi-turn conversations over the same five
repositories with per-turn chunk-level ground truth.
Each conversation contains 2--5 turns and at least one
coreference, one follow-up drill-down, or one topic-shift turn.
This set is used to evaluate session coherence, redundancy
suppression, and MRR under conversational conditions.

\subsection{Baselines}

\begin{itemize}[leftmargin=*]
  \item \textbf{Baseline}: Ollama local LLM, fixed 300-character
        sliding-window chunks, keyword (BM25) retrieval.
  \item \textbf{BM25}: Okapi BM25 ($k_1{=}2.0$, $b{=}0.75$),
        retrieval only, no LLM.
  \item \textbf{VanillaRAG}: BGE-small-en-v1.5 embeddings,
        cosine top-$k$ retrieval, no graph signals, no session
        state, Llama-3.1-8b generation.
  \item \textbf{\chatgit (Session\,+\,Intent)}: session memory
        and intent-adaptive retrieval without any graph
        components (no PageRank, no HITS, no call-context
        augmentation).
\end{itemize}

\subsection{Metrics}

\textbf{Setting A}: P@5 (precision at 5 retrieved chunks),
Recall@10, Accuracy (exact-match at position 1), and F1.

\textbf{Setting B}: MRR (mean reciprocal rank),
Recall@5, NDCG@5, P@1, and cross-turn redundancy (fraction of
turns where the top-ranked chunk was identical to a prior turn's
top-ranked chunk).

Statistical significance: Wilcoxon signed-rank test, bootstrap
confidence intervals ($n{=}3{,}000$ resamples), Bonferroni
correction.

% ══════════════════════════════════════════════════════════
\section{Results}
\label{sec:results}

\subsection{Setting A: System Evolution}

Table~\ref{tab:evolution} traces the improvement from the
initial baseline to the full system.  The most significant jump
occurs between the baseline and \chatgit v1, which replaces
keyword retrieval with dense embeddings, structured AST
chunking, and cross-encoder reranking.  The full system adds
session memory, intent routing, graph-structural signals, and
call-context augmentation on top of v1.

\begin{table}[t]
\centering
\caption{Setting A: system evolution (200 queries, 5 repos).
$\Delta$: full system vs.\ Baseline.}
\label{tab:evolution}
\small
\begin{tabular}{lcccc}
\toprule
System & P@5 & R@10 & Acc. & F1 \\
\midrule
\rowcolor{gray!10}
Baseline (BM25 + fixed chunks) & 0.48 & 0.55 & 0.62 & 0.60 \\
\chatgit v1 (BGE + PageRank + CE) & 0.78 & 0.86 & 0.89 & 0.88 \\
\rowcolor{cIndigo!10}
\textbf{\chatgit (full)} & \textbf{0.78} & \textbf{0.86}
  & \textbf{0.89} & \textbf{0.88} \\
\midrule
$\Delta$ (full vs.\ Baseline) & $+$63\% & $+$56\% & $+$44\% & $+$47\% \\
\bottomrule
\end{tabular}
\end{table}

\subsection{Setting B: Multi-Turn Evaluation}

Table~\ref{tab:multiturn} shows results on the 18-conversation
multi-turn set.  \chatgit achieves MRR 0.389 versus 0.104 for
BM25 ($p{=}0.013$, Cohen's $d{=}0.847$, large effect).
Cross-turn redundancy drops to zero for both \chatgit
configurations, confirming that session-memory redundancy
penalties reliably suppress re-retrieval of already-seen chunks.

\begin{table}[t]
\centering
\caption{Setting B: multi-turn evaluation ($n{=}18$ conversations).
\textbf{Bold}: best. $^{**}p{<}0.02$.}
\label{tab:multiturn}
\small
\begin{tabular}{lccccc}
\toprule
System & MRR & R@5 & NDCG@5 & P@1 & Redund. \\
\midrule
BM25
  & 0.104$_{\pm.083}$ & 0.167 & 0.099 & 0.000 & 9.4\% \\
VanillaRAG
  & 0.318$_{\pm.179}$ & 0.343 & 0.278 & 0.222 & 12.2\% \\
\chatgit (Sess.+Intent)
  & 0.347$_{\pm.195}$ & 0.444 & 0.346 & 0.278 & \textbf{0.0\%} \\
\rowcolor{cIndigo!10}
\textbf{\chatgit (full)}
  & \textbf{0.389}$_{\pm.215}$
  & \textbf{0.444} & \textbf{0.371} & \textbf{0.333}
  & \textbf{0.0\%} \\
\midrule
\multicolumn{6}{l}{\small vs.\ BM25: $\Delta{=}{+}0.285$,
  $p{=}0.013^{**}$, $d{=}0.847$ (large effect)} \\
\multicolumn{6}{l}{\small vs.\ VanillaRAG: $\Delta{=}{+}0.071$,
  $p{=}0.477$ (not significant)} \\
\bottomrule
\end{tabular}
\end{table}

\subsection{Ablation Study}

Table~\ref{tab:ablation} evaluates each component's contribution
by incrementally adding it over the Session+Intent base.
Bidirectional call-context augmentation provides the largest
single improvement (+0.042 MRR, +8\% relative).
Volatility weighting and PageRank show neutral or slightly
negative effect on this benchmark because test repositories
were shallow-cloned, eliminating the multi-commit history
needed for volatility signals and producing sparse call graphs
(discussed in \S\ref{sec:disc}).

\begin{table}[t]
\centering
\caption{Ablation study over the Session\,+\,Intent base.
$\dagger$: shallow-clone limitation; see \S\ref{sec:disc}.}
\label{tab:ablation}
\small
\begin{tabular}{lcccc}
\toprule
Configuration & MRR & R@5 & NDCG@5 & $\Delta$MRR \\
\midrule
VanillaRAG                      & 0.318 & 0.343 & 0.278 & $-0.029$ \\
Session\,+\,Intent              & 0.347 & 0.444 & 0.346 & --- \\
\quad +\,Volatility weight      & 0.340 & 0.444 & 0.339 & $-0.007^{\dagger}$ \\
\quad +\,Hybrid PageRank        & 0.340 & 0.444 & 0.342 & $-0.007^{\dagger}$ \\
\quad +\,Call-context augment.  & \textbf{0.389}
                                & \textbf{0.444}
                                & \textbf{0.371}
                                & \textbf{+0.042} \\
\chatgit (full)                 & \textbf{0.389}
                                & \textbf{0.444}
                                & \textbf{0.371}
                                & \textbf{+0.042} \\
\bottomrule
\end{tabular}
\end{table}

\subsection{Per-Intent and Per-Repository Analysis}

Table~\ref{tab:per_intent} shows per-intent MRR.  The full
system gains most on EXPLAIN (+17.7\,pp) and DEBUG (+9.0\,pp),
where call-context augmentation is most active.
SUMMARIZE collapses to 0.000 due to a ground-truth artefact
described in \S\ref{sec:disc}.
LOCATE is already strong under VanillaRAG (0.500); the
additional signals provide no measurable further gain on this
small set.

\begin{table}[t]
\centering
\caption{Per-intent MRR ($n{=}18$).
$\Delta$: \chatgit (full) minus VanillaRAG.}
\label{tab:per_intent}
\small
\begin{tabular}{lcccc}
\toprule
System & LOCATE & EXPLAIN & DEBUG & SUMMARIZE \\
\midrule
BM25              & 0.167 & 0.013 & 0.250 & 0.056 \\
VanillaRAG        & 0.500 & 0.229 & 0.348 & 0.250 \\
\chatgit (full)   & \textbf{0.500} & \textbf{0.406}
                  & \textbf{0.438} & 0.000 \\
$\Delta$          & 0.000 & \textbf{+0.177}
                  & \textbf{+0.090} & $-0.250$ \\
\bottomrule
\end{tabular}
\end{table}

Table~\ref{tab:per_repo} shows per-repository MRR.  FastAPI
yields the largest gain (+0.667) due to its decorator-heavy
architecture producing clean, traversable call edges.  Click
performs below BM25; its flat argument-parser design generates
sparse call graph edges that provide little retrieval signal.

\begin{table}[t]
\centering
\caption{Per-repository MRR, full system vs.\ BM25.}
\label{tab:per_repo}
\small
\begin{tabular}{lccc}
\toprule
Repository & \chatgit (full) & BM25 & $\Delta$ \\
\midrule
fastapi  & \textbf{0.667} & 0.000 & +0.667 \\
flask    & 0.542          & 0.139 & +0.403 \\
requests & 0.500          & 0.278 & +0.222 \\
celery   & 0.083          & 0.000 & +0.083 \\
click    & 0.000          & 0.070 & $-$0.070 \\
\midrule
\multicolumn{4}{l}{\small Wilcoxon vs.\ BM25: $p{=}0.013$,
  $d{=}0.847$} \\
\bottomrule
\end{tabular}
\end{table}

% ══════════════════════════════════════════════════════════
\section{Discussion}
\label{sec:disc}

\paragraph{Call-context augmentation is the dominant driver.}
The bidirectional call-neighbourhood (+4.2\,pp MRR) directly
addresses the most frequent retrieval failure mode: a function
chunk in isolation is insufficient for explain and debug
questions.  The proportional boost in Eq.~\eqref{eq:callctx}
ties neighbour weight to seed confidence, avoiding the
degradation observed with a fixed boost.

\paragraph{Volatility and PageRank on shallow clones.}
Both signals require a meaningful commit graph.
Shallow-cloned repositories with only a single commit provide
no per-file frequency or recency variance; every file receives
an identical volatility score.  Call graphs derived from
shallow clones may also miss import edges from files not
present in the working tree.  On full-history repositories,
these signals are expected to differentiate architectural layers
more strongly.

\paragraph{HITS separates roles that PageRank conflates.}
PageRank assigns a single global centrality score that rewards
files with many in-links.  HITS separates orchestrators
(high out-relevance, hub score) from core utilities (high
in-relevance, authority score), enabling intent-specific
boosting.  On repositories with clear architectural layering
(FastAPI, Flask), routing SUMMARIZE queries to hub nodes
and LOCATE/EXPLAIN queries to authority nodes produces
more coherent retrieved sets.

\paragraph{SUMMARIZE: a ground-truth gap.}
MRR collapses to 0.000 for SUMMARIZE because module-summary
chunks were excluded from the evaluation ground-truth set to
reduce annotation effort, making the correct answer unreachable
for any system.  The live pipeline generates and surfaces these
chunks correctly; this is a limitation of our small evaluation
construction, not of the retrieval approach.

\paragraph{Limitations.}
(1)~The multi-turn evaluation set contains only 18 conversations.
We report bootstrap CIs and effect sizes to quantify uncertainty,
but the results should be treated as indicative rather than
definitive.
(2)~The intent classifier is keyword-based and will misclassify
boundary-case queries; a trained classifier is expected to
improve both precision and recall on hybrid intents.
(3)~Shallow-clone test repositories suppress two of the
structural signals; a full evaluation on production-history
repositories remains future work.

% ══════════════════════════════════════════════════════════
\section{Conclusion}
\label{sec:conclusion}

We presented \chatgit, a multi-turn conversational code QA
system that addresses cross-turn redundancy, structurally-blind
retrieval, and fixed granularity through a cohesive pipeline of
complementary components.  Starting from a keyword-search
baseline (F1\,=\,0.60), the system reaches \textbf{F1\,=\,0.88}
on a 200-query single-turn suite (+47\% absolute) and
\textbf{MRR\,=\,0.389} on an 18-conversation multi-turn set
($3.7\times$ BM25, $p{=}0.013$, Cohen's $d{=}0.847$), with
zero cross-turn redundancy throughout.

The most impactful single component is bidirectional
call-context augmentation, which addresses the structural
incompleteness of isolated function chunks for explanation and
debugging tasks.  Session-aware retrieval memory eliminates
cross-turn repetition without any additional LLM calls.

\paragraph{Future work.}
Neural intent classification trained on a larger annotated
code QA corpus; hybrid sparse-dense retrieval; extending AST
parsing to Rust, C\#, and Kotlin; a controlled user study with
professional developers; and evaluation on full-history
production repositories to fully characterise the volatility
and PageRank components.

\paragraph{Acknowledgements.}
This project was implemented by Anish Gupta, Prakhar Sethi,
and Ritwik Bhattacharya under the supervision of
Dr.~Sonia Khetarpaul. We thank Dr.~Khetarpaul for her
guidance throughout.

% ══════════════════════════════════════════════════════════
\begin{thebibliography}{25}

\bibitem{githubcopilot}
GitHub: GitHub Copilot Chat.
\url{https://docs.github.com/en/copilot/github-copilot-chat} (2024)

\bibitem{cody}
Sourcegraph: Cody --- AI Coding Assistant.
\url{https://sourcegraph.com/cody} (2024)

\bibitem{cursor}
Anysphere Inc.: Cursor --- The AI Code Editor.
\url{https://www.cursor.com} (2024)

\bibitem{aider}
Gauthier, P.: Aider: AI Pair Programming in your Terminal.
\url{https://aider.chat} (2024)

\bibitem{continuedev}
Continue.dev: Open-Source Autopilot for Software Development.
\url{https://continue.dev} (2024)

\bibitem{lewis2021rag}
Lewis, P., et~al.: Retrieval-Augmented Generation for
Knowledge-Intensive NLP Tasks. arXiv:2005.11401 (2021)

\bibitem{petroni2024irrag}
Petroni, F., Siciliano, F., Silvestri, F., Trappolini, G.:
IR-RAG @ SIGIR24: Information Retrieval's Role in RAG Systems.
In: Proc.\ 47th ACM SIGIR, pp.\ 3036--3039 (2024)

\bibitem{guo2024lightrag}
Guo, Z., et~al.: LightRAG: Simple and Fast
Retrieval-Augmented Generation. arXiv:2410.05779 (2025)

\bibitem{robertson2009bm25}
Robertson, S., Zaragoza, H.: The Probabilistic Relevance
Framework: BM25 and Beyond.
Found.\ Trends IR \textbf{3}(4), 333--389 (2009)

\bibitem{feng2020codebert}
Feng, Z., et~al.: CodeBERT: A Pre-Trained Model for Programming
and Natural Languages. In: Findings of EMNLP (2020)

\bibitem{guo2021graphcodebert}
Guo, D., et~al.: GraphCodeBERT: Pre-Training Code Representations
with Data Flow. In: ICLR (2021)

\bibitem{zhang2023repocoder}
Zhang, F., et~al.: RepoCoder: Repository-Level Code Completion
Through Iterative Retrieval and Generation. In: NeurIPS (2023)

\bibitem{liu2024repoagent}
Liu, B., et~al.: RepoAgent: An LLM-Powered Framework for
Repository-level Code Documentation Generation.
arXiv:2402.16667 (2024)

\bibitem{kleinberg1999hits}
Kleinberg, J.M.: Authoritative Sources in a Hyperlinked
Environment. J.\ ACM \textbf{46}(5), 604--632 (1999)

\bibitem{brin1998pagerank}
Brin, S., Page, L.: The Anatomy of a Large-Scale Hypertextual
Web Search Engine. Comput.\ Netw.\ ISDN Syst. (1998)

\bibitem{velickovic2018gat}
Veli\v{c}kovi\'{c}, P., et~al.: Graph Attention Networks.
In: ICLR (2018)

\bibitem{reimers2019sentencebert}
Reimers, N., Gurevych, I.: Sentence-BERT: Sentence Embeddings
using Siamese BERT-Networks. In: EMNLP (2019)

\bibitem{nogueira2019passage}
Nogueira, R., Cho, K.: Passage Re-ranking with BERT.
arXiv:1901.04085 (2019)

\bibitem{reddy2019coqa}
Reddy, S., et~al.: CoQA: A Conversational Question Answering
Challenge. TACL \textbf{7}, 249--266 (2019)

\bibitem{choi2018quac}
Choi, E., et~al.: QuAC: Question Answering in Context.
In: EMNLP (2018)

\bibitem{jimenez2024swebench}
Jimenez, C.E., et~al.: SWE-bench: Can Language Models Resolve
Real-World GitHub Issues? In: ICLR (2024)

\bibitem{bavota2015changeadvisor}
Bavota, G., et~al.: How the Apache Community Upgrades
Dependencies. In: MSR (2015)

\bibitem{hindi2025rag}
Hindi, M., et~al.: Enhancing Precision and Interpretability of
RAG in Legal Technology. IEEE Access (2025)

\bibitem{abootorabi2025multimodal}
Abootorabi, M.M., et~al.: Ask in Any Modality: A Comprehensive
Survey on Multimodal RAG. arXiv:2502.08826 (2025)

\end{thebibliography}

\end{document}
